\chapter*{Resumen}
\addcontentsline{toc}{chapter}{Resumen}

Este proyecto integrador aborda el diseño y desarrollo de un sistema de alerta temprana de siniestros viales basado en inteligencia artificial, en el marco del convenio entre el Laboratorio de Investigación y Desarrollo de Software e Inteligencia Artificial (LIDeSIA) de la Facultad de Ciencias Exactas, Físicas y Naturales de la Universidad Nacional de Córdoba (FCEFyN-UNC) y el Ministerio Público Fiscal de la Provincia de Córdoba. El problema surge de la necesidad de la Unidad Judicial de Accidentología Vial de asignar recursos especializados, físicos y humanos (oficiales de justicia, peritos, personal policial y sanitario) ante eventos cuya ocurrencia es muy difícil anticipar, lo que obliga al organismo a sobredimensionar su operación.

El proyecto se gestionó con una metodología ágil híbrida: una etapa de descubrimiento con iteraciones semanales junto al cliente, al estilo de Scrum, y una etapa de ejecución organizada con un tablero Kanban. A partir de 4320 siniestros georreferenciados registrados entre 2022 y 2024, se reestructuró el conjunto de datos desde un formato de tabla de hechos hacia una formulación autorregresiva supervisada, agregando la información en bloques discretos de seis horas y ventanas deslizantes. Se experimentó con distintos modelos de inteligencia artificial (LightGBM, perceptrón multicapa, modelos georreferenciados y el pipeline automatizado de AutoGluon-TimeSeries) y se adoptó como modelo de línea base un XGBoost con regresión de Poisson y una capa de decisión posterior, priorizando la explicabilidad, en línea con la norma ISO/IEC 42001.

En el conjunto de prueba, el \qty{71.1}{\percent} de los bloques registró al menos una víctima, por lo que la situación actual del organismo equivale a una alerta permanente. Frente a esa referencia, el modelo reconoció el \qty{38.5}{\percent} de los bloques sin víctimas manteniendo una sensibilidad del \qty{87.4}{\percent} en los bloques críticos, y redujo el error cuadrático un \qty{6}{\percent} respecto de una predicción constante. Se trata de una señal predictiva modesta. Su valor principal es establecer una línea base medible y validar la factibilidad del enfoque, no constituir todavía una herramienta de planificación. El análisis de importancia de variables mostró que el riesgo se explica principalmente por el comportamiento reciente de la serie y la franja horaria, mientras que las variables climáticas resultaron marginales. El trabajo incluye la gestión de riesgos y la evaluación de impacto conforme a ISO/IEC 42001 y 23894, alcanzando un nivel de madurez tecnológica TRL~3 (ISO 16290:2013), y propone un plan de implementación de diez meses hacia TRL~7, con su cronograma, los perfiles necesarios e indicadores para evaluar su impacto.

\medskip
\noindent\textbf{Palabras clave:} alerta temprana, siniestralidad vial, aprendizaje automático, series temporales, XGBoost, ISO/IEC 42001, gestión de producto, metodologías ágiles.
